\section{Result 2: fault-specific corrective escalation}
\label{sec:escalation}

The attribution result above is a JAM verdict-layer statement. We now extend the ELVES branching model itself.

\subsection{Effective honest correction}

ELVES's honest offspring mean is
\[
\lambda=(1-\gamma)F.
\]
That expression treats every non-adversarial auditor as capable of identifying the invalid report. Under fault \(x\), only the independently correct class \(C_x\) supplies useful corrective offspring. We therefore replace the honest share by
\[
h_x=\frac{C_x}{N}=(1-\gamma)(1-f_x).
\]
The fault-specific reproduction mean is
\[
\boxed{
\lambda_x=Fh_x=F(1-\gamma)(1-f_x).
}
\]
This is an extension of ELVES, not a result stated in \citet{burdges2024}.

The branching process remains supercritical only if
\[
\lambda_x>1,
\]
so
\[
\boxed{
 f_x< f_{\rm crit}
 =1-\frac{1}{(1-\gamma)F}.
}
\]

For \(F=2\):
\begin{center}
\begin{tabular}{ccc}
\toprule
Adversarial share \(\gamma\) & \(f_{\rm crit}\) & Interpretation \\
\midrule
\(1/3\) & 0.250 & 25\% of honest validators \\
0.20 & 0.375 & 37.5\% of honest validators \\
0.05 & 0.474 & 47.4\% of honest validators \\
\bottomrule
\end{tabular}
\end{center}

The concentration boundary should be interpreted as a \emph{fault-domain} limit, not a client-count rule. Client deployment share is an observable proxy only when the client is the relevant common failure domain.

\subsection{Target-specific margin}

Structural supercriticality is weaker than satisfying a quantitative soundness target. Let \(q_x\) solve
\[
q_x=\exp[-\lambda_x(1-q_x)]
\]
for \(\lambda_x>1\). Then the extended pessimistic bound is
\[
P_x\le q_x^{s_0/F}.
\]

ELVES Corollary 1 provides an economically derived target of the form
\[
\varepsilon_* = \frac{\nu}{N+\nu},
\]
when collateral is parameterized by \(\nu\) in the stated expected-profit argument. Inverting the branching bound for target \(\varepsilon\) gives
\[
q_* = \varepsilon^{F/s_0},
\qquad
\lambda_*=-\frac{\ln q_*}{1-q_*},
\]
and therefore
\[
\boxed{
 f_{\max}
 =1-\frac{\lambda_*}{(1-\gamma)F},
}
\]
provided the baseline \(f_x=0\) already meets the target.

For \(N=1023\), \(s_0=30\), \(F=2\), and illustrative \(\nu=0.05\), the target is not certified at the \(\gamma\to1/3\) boundary even at \(f_x=0\). At \(\gamma=0.20\), the bound certifies only about \(f_x\le0.146\); at \(\gamma=0.05\), about \(f_x\le0.280\). These numbers are certificate margins of the extended pessimistic model, not exact protocol frontiers.

\subsection{The tranche factor trades concentration tolerance against escalation load}

Increasing \(F\) raises \(\lambda_x\) and therefore tolerates a larger correlated-fault domain. But the ELVES static committee-size bound under a no-show/fault share \(u\) is
\[
 s_f\le\frac{s_0}{1-uF}.
\]
Thus the same parameter reduces distance to no-show explosion.

For \(s_0=30\) and \(u=0.30\):

\begin{table}[htbp]
\centering
\caption{Illustrative two-sided effect of the tranche factor. The client/fault thresholds are from the correlated-honest extension; the committee bound is the ELVES static form evaluated at \(u=0.30\).}
\label{tab:Ftradeoff}
\begin{tabular}{rrrr}
\toprule
\(F\) & \(f_{\rm crit}\), \(\gamma=1/3\) & \(f_{\rm crit}\), \(\gamma=0.20\) & \(s_f\) bound \\
\midrule
2.00 & 25.0\% & 37.5\% & 75 \\
2.0794 & 27.9\% & 39.9\% & 79.7 \\
2.50 & 40.0\% & 50.0\% & 120 \\
2.90 & 48.3\% & 56.9\% & 230.8 \\
\bottomrule
\end{tabular}
\end{table}

With \(s_0=30\) and \(\gamma=1/3\), the pessimistic acceptance bound is approximately \(1.13\times10^{-4}\) at \(F=2\) and \(4.54\times10^{-5}\) at ELVES's reported continuous optimum \(F\approx2.0794\), a ratio of about 2.49. This is a comparison of one modeled bound, not a statement that JAM is “2.5 times less secure.”

The design implication is simply that \(F\) sits between two margins:
\[
\boxed{
\text{corrective reproduction}
\quad\text{and}\quad
\text{failure/no-show reproduction}.
}
\]
If correlated honest failure is admitted to the threat model, the objective used to select \(F\) changes.
